\section{Related Work}

\paragraph{MAST taxonomy.} The Multi-Agent System Trace (MAST) taxonomy
\cite{mast2024} catalogs 14 failure modes across three coarse categories:
specification issues (1.x), inter-agent misalignment (2.x), and task verification
issues (3.x). We use the \texttt{agentdash} reference annotator to label each trial.

\paragraph{Scaling agent systems.} \cite{scaling_agents_2025} sweep aggregate failure
rates across multi-agent benchmarks and report 41--86\% failures depending on
benchmark and stack, but their study does not isolate protocol or topology as a
controlled factor; we treat their range as the prior baseline against which our
per-cell FC2 rates should be compared.

\paragraph{Agent communication protocols.} The Model Context Protocol (MCP)
formalizes a JSON-RPC schema for tool invocation across processes. Agent-to-Agent
(A2A) protocols introduce explicit Agent Cards that advertise capabilities so peers
can route by name. Native function-call multi-agent stacks (e.g.\ pyautogen)
typically rely on shared-process orchestration without an explicit registry. The
prevailing question for practitioners is whether these protocol differences
\emph{matter} for failure modes once the topology is held fixed; this paper tests
that question empirically.

\paragraph{Topology effects.} Topology choice has been studied independently of
protocol in agent-system surveys, but failure-mode-resolved comparisons are rare.
Our $3 \times 3$ design is, to our knowledge, the first that jointly varies both
axes under a single MAST annotator and a fixed task suite.

\paragraph{ProtocolBench and SEMAP.} \cite{protocolbench2025} compares
A2A, ACP, ANP, and Agora on aggregate success rate and latency; it
\emph{excludes} MCP and does not vary topology. \cite{semap2025} studies
self-evolving A2A-only stacks under a single aggregate failure rate. We
add MCP, add the topology axis, and replace the success-rate / aggregate
metric with the per-MAST-category breakdown — turning what was a
benchmark into a diagnosis tool.
